\appendices
\section{Detailed Seed Breakdown on OpTC}
Table~\ref{tab:seed_breakdown} reports the per-seed breakdown of classification AUROC and empirical False Positive Rate (FPR) for ZeroCausal (Tuned SCM) and the Isolation Forest baseline across 10 random seeds on the OpTC dataset.

\begin{table}[htbp]
\centering
\caption{Detailed Seed Breakdown on OpTC Dataset}
\begin{tabular}{cccc}
\toprule
Seed & ZeroCausal AUC & ZeroCausal FPR & Isolation Forest AUC \\
\midrule
42 & 0.8188 & 4.21\% & 0.6117 \\
43 & 0.6359 & 4.24\% & 0.6584 \\
44 & 0.6770 & 4.23\% & 0.6700 \\
45 & 0.6971 & 3.53\% & 0.6703 \\
46 & 0.7849 & 3.97\% & 0.6952 \\
47 & 0.7714 & 4.44\% & 0.6826 \\
48 & 0.7924 & 4.25\% & 0.6919 \\
49 & 0.7088 & 4.43\% & 0.6987 \\
50 & 0.5973 & 5.10\% & 0.6745 \\
51 & 0.7876 & 2.43\% & 0.6926 \\
\midrule
Mean & $0.7271 \pm 0.0711$ & $4.09\% \pm 0.66\%$ & $0.6746 \pm 0.0244$ \\
Median & 0.7401 & 4.23\% & 0.6786 \\
\bottomrule
\end{tabular}
\label{tab:seed_breakdown}
\end{table}

\section{Artifact Appendix}

This appendix provides instructions for reviewers and researchers to set up the environment, run the evaluation pipelines, and reproduce all results and figures presented in the paper.

\subsection{Abstract}
The artifact consists of the core Python implementation of \textbf{ZeroCausal}, evaluation scripts for five benchmarks (DARPA OpTC, simulated DARPA TC3, simulated NODLINK, BETH, and StreamSpot), configurations, and plotting tools. The evaluations run online causal discovery (PCMCI with ParCorr), fit Structural Causal Models (SCMs), monitor anomalies using the Causal Anomaly Score (CAS), and perform online conformal calibration. Using a standard Linux host, all benchmarks and figures can be generated in under 5 minutes (excluding the one-time high-dimensional PCMCI discovery on BETH).

\subsection{Artifact Meta-Information}
\textbf{Algorithm}: PCMCI causal discovery, SCM residual regression, conformal calibration, adaptive window change-point detection.
\textbf{Program}: Python 3.9--3.11 on Linux/macOS/WSL2. Default seed: \texttt{42}.
\textbf{Execution Time}: $\sim$28\,s for OpTC (streaming only; PCMCI baseline $\sim$5\,min, run once); $\sim$2--3\,s each for TC3 and NODLINK.
\textbf{Output}: JSON summary metrics, CSV step logs, PNG figures.

\subsection{Requirements}
\textbf{Hardware}: Standard x86/ARM64, $\geq$2 cores, $\geq$4\,GB RAM, $\geq$500\,MB storage.
\textbf{Software}: Python 3.9--3.11; dependencies listed in \texttt{requirements.txt}.

\subsection{Installation \& Environment Setup}
To run the artifact, first clone the repository and establish the virtual environment:
\begin{lstlisting}
python3 -m venv .venv_zc
source .venv_zc/bin/activate
pip install --upgrade pip
pip install -r requirements.txt
\end{lstlisting}

\subsection{Reproduction of Experimental Results}

\subsubsection{Run the Evaluation Pipeline}
\textbf{OpTC (seed 42):}
\begin{lstlisting}
python3 05_evaluate_zerocausal.py \
  --std-floor 1.0 --seed 42 --run-name optc_run
python3 run_all_seeds.py  # 10-seed aggregate
\end{lstlisting}

\textbf{TC3 (with concept-drift simulation):}
\begin{lstlisting}
python3 09_evaluate_additional_datasets.py \
  --dataset tc3 --baseline --simulate-drift \
  --std-floor 1.0 --run-name tc3_trace_default
\end{lstlisting}

\textbf{NODLINK:}
\begin{lstlisting}
python3 09_evaluate_additional_datasets.py \
  --dataset nodlink --baseline \
  --std-floor 1.0 --run-name nodlink_default
\end{lstlisting}

\textbf{BETH:}
\begin{lstlisting}
python3 09_evaluate_additional_datasets.py \
  --dataset beth --baseline \
  --std-floor 1.0 --run-name beth_default
\end{lstlisting}

\textbf{StreamSpot:}
\begin{lstlisting}
python3 09_evaluate_additional_datasets.py \
  --dataset streamspot --baseline \
  --std-floor 1.0 --run-name streamspot_default
\end{lstlisting}

\subsubsection{Generate and Replicate the Figures}
Run the plotting scripts to compile the vector figures.
\begin{lstlisting}
python3 generate_architecture_diagram.py
python3 10_plot_comparisons.py
python3 11_sensitivity_analysis.py
\end{lstlisting}
This generates \texttt{zerocausal\_architecture.png}, \texttt{benchmark\_comparison\_roc.png}, \texttt{threshold\_adaptation\_learning.png}, and \texttt{noise\_sensitivity\_analysis.png} in the \texttt{results/final/} directory.

\subsection{Code Availability}
An anonymised version of the ZeroCausal codebase is available for reviewer access at: \url{https://anonymous.4open.science/r/ZeroCausal-XXXX}.
